% ESTRUTURA INICIAL DE DISSERTAÇÃO - EXPERIMENTAL.
% Não submeter como versão final sem adaptar ao manual da universidade.
\documentclass[12pt,oneside,a4paper,brazil]{abntex2}
\usepackage[T1]{fontenc}
\usepackage[utf8]{inputenc}
\usepackage[brazil]{babel}
\usepackage{lmodern}
\usepackage{graphicx,amsmath,amssymb,booktabs}
\usepackage{hyperref,url}
\usepackage[alf]{abntex2cite}
\providecommand{\citep}[1]{\cite{#1}}
\providecommand{\citet}[1]{\citeonline{#1}}
\providecommand{\citealp}[1]{\cite{#1}}
\providecommand{\citealt}[1]{\citeonline{#1}}
\providecommand{\citeauthor}[1]{\citeonline{#1}}
\providecommand{\citeyear}[1]{\cite{#1}}
\titulo{Título provisório da dissertação}
\autor{Pesquisador(a) — informar nome}
\orientador{Professor(a) orientador(a) — informar}
\instituicao{Instituição e Programa de Pós-Graduação — informar}
\tipotrabalho{Dissertação (rascunho)}
\preambulo{Estrutura preliminar para organização de uma pesquisa de mestrado, pendente de adequação às orientações institucionais.}
\local{Brasil}
\data{\the\year}
\begin{document}
\selectlanguage{brazil}
\imprimircapa
\begin{titlingpage}
\centering
{\large\imprimirinstituicao\par}
\vspace*{3cm}
{\large\imprimirautor\par}
\vspace*{2cm}
{\Large\bfseries\imprimirtitulo\par}
\vfill
{\imprimirpreambulo\par}
\vspace*{1cm}
{\imprimirorientador\par}
\vfill
{\imprimirlocal, \imprimirdata\par}
\end{titlingpage}
\begin{resumo}
Este rascunho apresenta a proposta de estudo. Não apresente dados empíricos que não existam.
\vspace{\onelineskip}
\noindent\textbf{Palavras-chave:} pesquisa; metodologia; revisão bibliográfica.
\end{resumo}
\tableofcontents*
\textual
\chapter{Introdução}
\label{sec:intro}
%%%%%%%%% INTRODUÇÃO %%%%%%%%%
Apresente problema, perguntas e objetivos.
%%%%%%%%% INTRODUÇÃO %%%%%%%%%
\chapter{Referencial teórico}
\label{sec:related_work}
%%%%%%%%% REVISÃO %%%%%%%%%
Discussão bibliográfica sustentada por fontes verificáveis.
%%%%%%%%% REVISÃO %%%%%%%%%
\chapter{Procedimentos metodológicos}
\label{sec:method}
%%%%%%%%% MÉTODO %%%%%%%%%
Descreva os procedimentos, o corpus, o tratamento e as limitações.
%%%%%%%%% MÉTODO %%%%%%%%%
\chapter{Análise e discussão}
\label{sec:experiments}
%%%%%%%%% RESULTADOS %%%%%%%%%
Se não houver coleta/análise concluída, declare as limitações; não invente dados.
%%%%%%%%% RESULTADOS %%%%%%%%%
\chapter{Considerações finais}
\label{sec:conclusion}
%%%%%%%%% CONCLUSÃO %%%%%%%%%
Indique síntese e trabalhos futuros sem alegar resultados inexistentes.
%%%%%%%%% CONCLUSÃO %%%%%%%%%
\postextual
\bibliography{references}
\end{document}
